%! Matrices as Functions — Unit 4, Lesson 4.13 — Homework (Answer Key)
\documentclass[10pt]{article}
\usepackage{precalculus-article}
\usepackage{precalculus-key}   % pulls in -boxes; do NOT also load -boxes
\newcommand{\TallMath}[1]{$\displaystyle #1\rule[-1.4em]{0pt}{3.2em}$}

\begin{document}
\pageheader{Unit 4, Lesson 4.13}{Homework --- Answer Key}
\namedateperiod
\vspace{0.10in}

\begin{notesbox}
\textbf{1.} A transformation maps $\hat{\imath}\to\langle3,1\rangle$ and
$\hat{\jmath}\to\langle-2,4\rangle$.

\vspace{0.06in}
\textbf{(a)} Write the matrix $A$. \hspace{1em}
$A = $ \ans{$\begin{bmatrix}3&-2\\1&4\end{bmatrix}$}

\vspace{0.12in}
\textbf{(b)} Compute $L\bigl(\langle2,-1\rangle\bigr)$. Show your work.

\vspace{0.06in}
\ans{$\begin{bmatrix}3&-2\\1&4\end{bmatrix}\begin{bmatrix}2\\-1\end{bmatrix}
= \begin{bmatrix}3(2)+(-2)(-1)\\1(2)+4(-1)\end{bmatrix}
= \begin{bmatrix}8\\-2\end{bmatrix}$}

\vspace{0.04in}
$L\bigl(\langle2,-1\rangle\bigr) = $ \ans{$\begin{bmatrix}8\\-2\end{bmatrix}$}
\end{notesbox}

\vspace{0.10in}

\begin{notesbox}
\textbf{2.} Use $\cos30°=\dfrac{\sqrt{3}}{2}$ and $\sin30°=\dfrac{1}{2}$.

\vspace{0.06in}
\textbf{(a)} Write the rotation matrix $R_{30°}$.

\vspace{0.06in}
$R_{30°} = $ \ans{$\begin{bmatrix}\dfrac{\sqrt{3}}{2}&-\dfrac{1}{2}\\[6pt]\dfrac{1}{2}&\dfrac{\sqrt{3}}{2}\end{bmatrix}$}

\vspace{0.12in}
\textbf{(b)} Compute $\det(R_{30°})$.

\vspace{0.06in}
$\det(R_{30°}) = $ \ans{1} \quad (\ans{$\cos^2 30°+\sin^2 30° = \tfrac{3}{4}+\tfrac{1}{4}=1$})

\vspace{0.12in}
\textbf{(c)} Apply $R_{30°}$ to $\langle2,0\rangle$. Show your work.

\vspace{0.06in}
\ans{$R_{30°}\begin{bmatrix}2\\0\end{bmatrix}
= \begin{bmatrix}\tfrac{\sqrt{3}}{2}(2)+(-\tfrac{1}{2})(0)\\[4pt]\tfrac{1}{2}(2)+\tfrac{\sqrt{3}}{2}(0)\end{bmatrix}
= \begin{bmatrix}\sqrt{3}\\1\end{bmatrix}$}

\vspace{0.04in}
$R_{30°}\langle2,0\rangle = $ \ans{$\begin{bmatrix}\sqrt{3}\\1\end{bmatrix}$}
\end{notesbox}

\vspace{0.10in}

\begin{notesbox}
\textbf{3.} Let $A=\begin{bmatrix}2&1\\0&3\end{bmatrix}$ and
$B=\begin{bmatrix}1&0\\-1&2\end{bmatrix}$.

\vspace{0.06in}
\textbf{(a)} Compute $BA$ (apply $A$ first, then $B$). Show your work.

\vspace{0.06in}
\ans{$BA = \begin{bmatrix}1&0\\-1&2\end{bmatrix}\begin{bmatrix}2&1\\0&3\end{bmatrix}
= \begin{bmatrix}1(2)+0(0)&1(1)+0(3)\\(-1)(2)+2(0)&(-1)(1)+2(3)\end{bmatrix}
= \begin{bmatrix}2&1\\-2&5\end{bmatrix}$}

\vspace{0.04in}
$BA = $ \ans{$\begin{bmatrix}2&1\\-2&5\end{bmatrix}$}

\vspace{0.10in}
\textbf{(b)} What is $|\det(BA)|$?

\vspace{0.04in}
$|\det(BA)| = $ \ans{$|2(5)-1(-2)| = |10+2| = 12$} \quad $= $ \ans{12}
\end{notesbox}

\vspace{0.10in}

\begin{notesbox}
\textbf{4.} Let $L(\vec{v})=A\vec{v}$ with $A=\begin{bmatrix}3&1\\-1&2\end{bmatrix}$.

\vspace{0.06in}
\textbf{(a)} Compute $\det A$. \hspace{1em} $\det A = $ \ans{$3(2)-(1)(-1) = 7$}

\vspace{0.10in}
\textbf{(b)} Find $A^{-1}$ using the $2\times2$ inverse formula.

\vspace{0.06in}
\ans{$A^{-1} = \dfrac{1}{7}\begin{bmatrix}2&-1\\1&3\end{bmatrix}$}

\vspace{0.04in}
$A^{-1} = $ \ans{$\dfrac{1}{7}\begin{bmatrix}2&-1\\1&3\end{bmatrix}$}

\vspace{0.10in}
\textbf{(c)} Apply $L^{-1}$ to $\langle7,1\rangle$.

\vspace{0.06in}
\ans{$A^{-1}\begin{bmatrix}7\\1\end{bmatrix}
= \dfrac{1}{7}\begin{bmatrix}2&-1\\1&3\end{bmatrix}\begin{bmatrix}7\\1\end{bmatrix}
= \dfrac{1}{7}\begin{bmatrix}14-1\\7+3\end{bmatrix}
= \dfrac{1}{7}\begin{bmatrix}13\\10\end{bmatrix}
= \begin{bmatrix}13/7\\10/7\end{bmatrix}$}

\vspace{0.04in}
$L^{-1}\bigl(\langle7,1\rangle\bigr) = $ \ans{$\left\langle\dfrac{13}{7},\,\dfrac{10}{7}\right\rangle$}
\end{notesbox}

\vspace{0.10in}

\begin{notesbox}
\textbf{5. AP-Style Multiple Choice.}
The matrix $\begin{bmatrix}\cos\theta&-\sin\theta\\\sin\theta&\cos\theta\end{bmatrix}$
always has determinant equal to:

\vspace{0.08in}
\begin{itemize}[leftmargin=2em, itemsep=4pt]
  \item[(A)] 0
  \item[(B)] $\theta$
  \item[(C)] 1
  \item[(D)] $\cos\theta$
\end{itemize}

\vspace{0.08in}
Answer: \ans{(C)} --- $\det = \cos^2\theta + \sin^2\theta = 1$ by the Pythagorean identity.
\end{notesbox}

\end{document}
